% !TeX root = ../../Book1.tex
% 纲要标题：## 第6章　直积、半直积与群扩张
% 纲要位置：计划纲要.md，第 275 行。

\chapter{直积、半直积与群扩张}
\label{b1:ch:06}
\BookChapterNumber{b1:ch:06}

\begin{chapterabstract}
知道两个群后，怎样构造一个以它们为组成部分的新群？直积使两个因子独立运算，半直积加入一个自同构作用，群扩张则允许陪集代表的乘法产生额外因子。本章建立直积、直和与半直积的构造及其映射性质，再用截面和因子集描述一般群扩张，并通过二面体群、仿射群、四元数群和循环群说明分裂与非分裂的区别。
\end{chapterabstract}

\begin{learninggoals}
\begin{itemize}
\item 区分任意直积、有限支集直积和交换群直和的泛性质。
\item 验证外半直积的群公理，运用内半直积判别计算具体群。
\item 从短正合群列识别核、商及截面，说明扩张等价要求保持哪些映射，证明分裂与半直积分解等价。
\item 计算简单扩张的因子集，辨认依赖截面选择的数据和不变的结构。
\end{itemize}
\end{learninggoals}

本章需要用到\chapref{b1:ch:02}的正规子群和商群、\chapref{b1:ch:03}的自同构与具体有限群，以及\chapref{b1:ch:04}的共轭作用。半直积本身不依赖 Sylow 理论；\chapref{b1:ch:05}的 \(pq\) 阶群分类在本章只作为已完成的算例重新解释。

\ProseParagraphBreak
初读时先掌握\secref{b1:sec:0601}和\secref{b1:sec:0602}的直积、直和与半直积，以及\cref{b1:thm:split-extension-semidirect}的分裂判别和\cref{b1:cor:abelian-kernel-action}的交换核诱导作用；再用\subsecref{b1:subsec:0604:01}和\subsecref{b1:subsec:0604:02}的具体扩张检验这些概念。\subsecref{b1:subsec:0603:03}允许核不交换，把截面、共轭和结合律合在一起描述一般扩张，可作为综合训练；继续学习下一章的正规列与可解群时，可以暂缓其中的一般因子集计算。\subsecref{b1:subsec:0604:03}末尾的上同调解释是后续学习的预告，不是本章证明的前提。

\ProseParagraphBreak
直积与扩张可参阅 Rotman 的《An Introduction to the Theory of Groups》\cite[Chapters 6--7]{Rotman1995TheoryGroups}。若要继续研究分裂扩张与交换核扩张的上同调分类，可在学习模与基本同调代数后阅读 Brown 的《Cohomology of Groups》\cite[Chapter IV, \S\S~2--3, pp. 87--96]{Brown1982CohomologyGroups}；非交换核扩张的概述见同章\cite[Chapter IV, \S~6, pp. 104--106]{Brown1982CohomologyGroups}。本章先用具体乘法说明截面、半直积和因子集，不以群上同调为前提。

% !TeX root = ../../Book1.tex
% 纲要标题：### 6.1 直积与直和
% 纲要位置：计划纲要.md，第 273 行。

\section{直积与直和}
\label{b1:sec:0601}

% TODO: 按以下纲要要点撰写本节。

% 纲要内容（仅作写作计划，尚非教材正文）：
%
% * 有限直积及其投影
% * 任意群族的直积与限制直积式直和
% * 阿贝尔群范畴中的直和
% * 非交换群的有限支集直积是直积型构造，不是群范畴中的余积；后者由第8.3节的自由积承担
%

\subsection{有限直积及其投影}
\label{b1:subsec:0601:01}

待撰写。

\subsection{任意直积与有限支集直积}
\label{b1:subsec:0601:02}

待撰写。

\subsection{阿贝尔群的直和}
\label{b1:subsec:0601:03}

待撰写。

\subsection{有限支集直积与群的余积}
\label{b1:subsec:0601:04}

待撰写。

% !TeX root = ../../Book1.tex
% 纲要标题：### 6.2 半直积
% 纲要位置：计划纲要.md，第 280 行。

\section{半直积}
\label{b1:sec:0602}

% TODO: 按以下纲要要点撰写本节。

% 纲要内容（仅作写作计划，尚非教材正文）：
%
% * 自同构作用与外半直积
% * 内半直积判别
% * 二面体群、仿射群与具体计算
%

\subsection{自同构作用与外半直积}
\label{b1:subsec:0602:01}

待撰写。

\subsection{内半直积判别}
\label{b1:subsec:0602:02}

待撰写。

\subsection{二面体群、仿射群与具体计算}
\label{b1:subsec:0602:03}

待撰写。

% !TeX root = ../../Book1.tex
% 纲要标题：### 6.3 群扩张
% 纲要位置：计划纲要.md，第 286 行。

\section{群扩张}
\label{b1:sec:0603}

% 本节正文按下列原纲要要点撰写。

% 纲要内容（仅作写作计划，尚非教材正文）：
%
% * 短正合群列与分裂扩张
% * 截面是集合映射还是群同态
% * 扩张数据与半直积的关系
%

\subsection{短正合群列与分裂扩张}
\label{b1:subsec:0603:01}

商群把正规子群压缩为单位；群扩张提出相反的问题：已知正规子群和商群，能恢复哪些群？直积和半直积给出两种答案，但并非所有扩张都属于这两类。

\begin{definition}\label{b1:def:group-extension}
一列群同态
\[
1\longrightarrow N\xrightarrow{\iota}G
\xrightarrow{\pi}H\longrightarrow1
\]
称为\term[duan-zhenghe-qunlie]{短正合群列}[short exact sequence of groups]，如果 \(\iota\) 单射、\(\pi\) 满射，且 \(\operatorname{im}\iota=\ker\pi\)。它也称为 \(H\) 由 \(N\) 的一个\term[qun-kuozhang]{群扩张}[group extension]。若存在群同态 \(s:H\rightarrow G\) 满足 \(\pi s=\operatorname{id}_H\)，称扩张\term[fenlie-kuozhang]{分裂}[split extension]。
\end{definition}
这里的 \(1\) 表示平凡群。“正合”在每个位置要求前一个映射的像等于后一个映射的核，因此两端的条件正好是单射与满射。通过 \(\iota\)，可将 \(N\) 识别为 \(G\) 的正规子群；此时 \(G/N\cong H\)。但扩张还记住了指定映射 \(\iota,\pi\)，比只写一个抽象同构 \(G/N\cong H\) 包含更多信息。

\begin{theorem}\label{b1:thm:split-extension-semidirect}
短正合群列分裂，当且仅当 \(G\) 含有子群 \(L\)，使 \(\pi|_L:L\rightarrow H\) 为同构。存在这样的 \(L\) 时，\(G\) 是 \(N\) 与 \(L\) 的内半直积；若选择了同态截面 \(s\)，则
\[
G\cong N\rtimes_\alpha H,\qquad
\alpha_h(n)=s(h)ns(h)^{-1},
\]
同构由 \((n,h)\mapsto ns(h)\) 给出。
\end{theorem}
\begin{proof}
若有同态截面 \(s\)，则由 \(\pi s=\operatorname{id}\) 可知 \(s\) 单射，像 \(L=s(H)\) 是子群，且 \(\pi|_L\) 与 \(s\) 互逆。对任意 \(g\in G\)，令 \(h=\pi(g)\)，则 \(gs(h)^{-1}\in\ker\pi=N\)，所以 \(G=NL\)。若 \(s(h)\in N\)，应用 \(\pi\) 得 \(h=1\)，故 \(N\cap L=\{1\}\)。内半直积判别给出所需同构。

反过来，若 \(\pi|_L\) 是同构，取它的逆再接子群嵌入，得到同态截面。特别地，外半直积投影 \(N\rtimes H\rightarrow H\) 总有同态截面 \(h\mapsto(1,h)\)。
\end{proof}

\subsection{集合截面与同态截面}
\label{b1:subsec:0603:02}

满射 \(\pi:G\rightarrow H\) 的\term[jiemian]{截面}[section]，作为集合映射，仅要求为每个 \(h\) 选择一个原像 \(s(h)\)。有限群时可逐个选择，一般情形可借助选择公理。本节总把单位的代表选为单位，即 \(s(1)=1\)。集合截面存在不等于群扩张分裂，后者要求这个选择还能保持乘法。

\ProseParagraphBreak
考虑加法群的满射
\[
\mathbb Z/4\mathbb Z\longrightarrow\mathbb Z/2\mathbb Z,
\qquad [a]_4\longmapsto[a]_2.
\]
选择 \(s([0]_2)=[0]_4\)、\(s([1]_2)=[1]_4\) 得到集合截面，但
\[
s([1]_2+[1]_2)=[0]_4,\qquad
s([1]_2)+s([1]_2)=[2]_4,
\]
两者不同。另一种可能的非零陪集代表为 \([3]_4\)，其加倍仍为 \([2]_4\)，所以也不能给出同态截面。这里核与商群都同构于 \(C_2\)，中间群却是 \(C_4\)，并非 \(C_2\times C_2\)。

\ProseParagraphBreak
对任意归一化集合截面 \(s\)，每个 \(g\in G\) 仍有唯一表达 \(g=ns(h)\)：先令 \(h=\pi(g)\)，再令 \(n=gs(h)^{-1}\in N\)。因此 \(G\) 的底集合总可用 \(N\times H\) 的坐标描述；困难在于乘法，而非集合的分解。与半直积相比，新出现的项正好记录截面未能保持乘法的差额。

\subsection{扩张数据与半直积}
\label{b1:subsec:0603:03}

给定归一化集合截面，定义
\[
\alpha_h(n)=s(h)ns(h)^{-1},\qquad
c(h,k)=s(h)s(k)s(hk)^{-1}\in N.
\]
每个 \(\alpha_h\) 都是 \(N\) 的自同构，但映射 \(h\mapsto\alpha_h\) 未必是同态。元素 \(c(h,k)\) 称为这个截面的\term[yinzi-ji]{因子集}[factor set]；它等于单位当且仅当截面对相应的一对元素保持乘法。

\begin{proposition}\label{b1:prop:extension-factor-data}
在坐标 \(g=ns(h)\) 下，扩张中的乘法为
\[
(n,h)(m,k)=(n\alpha_h(m)c(h,k),hk).
\]
这些数据满足 \(\alpha_1=\operatorname{id}\)、\(c(1,h)=c(h,1)=1\)，并满足
\begin{align}
\alpha_h\alpha_k
&=\operatorname{Inn}_{c(h,k)}\alpha_{hk},\label{b1:eq:extension-action-defect}\\
c(h,k)c(hk,\ell)
&=\alpha_h(c(k,\ell))c(h,k\ell),\label{b1:eq:extension-factor-identity}
\end{align}
其中 \(\operatorname{Inn}_u(n)=unu^{-1}\) 表示由 \(u\) 给出的\term[nei-zitonggou]{内自同构}[inner automorphism]。
\end{proposition}
\begin{proof}
先把 \(m\) 移过 \(s(h)\)，再用 \(s(h)s(k)=c(h,k)s(hk)\)，得到
\[
ns(h)ms(k)=n\alpha_h(m)c(h,k)s(hk),
\]
这给出乘法公式。由 \(s(h)s(k)=c(h,k)s(hk)\)，两边对 \(N\) 作共轭便得到\eqref{b1:eq:extension-action-defect}。另一方面，将 \(s(h)s(k)s(\ell)\) 分别按左右结合展开，左结合为
\[
c(h,k)c(hk,\ell)s(hk\ell),
\]
右结合为
\[
\alpha_h(c(k,\ell))c(h,k\ell)s(hk\ell).
\]
消去右端共同因子，就得到\eqref{b1:eq:extension-factor-identity}。归一化条件直接来自 \(s(1)=1\)。
\end{proof}
若把截面改为 \(s'(h)=b(h)s(h)\)，其中 \(b(h)\in N\)、\(b(1)=1\)，则直接代入定义给出
\[
\alpha'_h=\operatorname{Inn}_{b(h)}\alpha_h,\qquad
c'(h,k)=b(h)\alpha_h(b(k))c(h,k)b(hk)^{-1}.
\]
计算第二式时，唯一需要移过截面的因子是 \(b(k)\)，由
\(s(h)b(k)=\alpha_h(b(k))s(h)\) 处理。可见因子集依赖截面的选择，同一个扩张可以有不同的 \(c\)。扩张分裂当且仅当能选到使全部 \(c(h,k)=1\) 的截面；这时\eqref{b1:eq:extension-action-defect}说明 \(\alpha\) 成为真正的同态，乘法恢复为半直积。对于非交换的 \(N\)，不能删掉内自同构项并把任意 \(\alpha\) 都当作群作用。

% !TeX root = ../../Book1.tex
% 纲要标题：### 6.4 扩张的例子
% 纲要位置：计划纲要.md，第 292 行。

\section{扩张的例子}
\label{b1:sec:0604}

% TODO: 按以下纲要要点撰写本节。

% 纲要内容（仅作写作计划，尚非教材正文）：
%
% * 中心扩张与非分裂例子
% * 四元数群与二面体群的对照
% * 第四卷第二群上同调的分类解释预告
%
% ---
%

\subsection{中心扩张与非分裂例子}
\label{b1:subsec:0604:01}

待撰写。

\subsection{四元数群与二面体群}
\label{b1:subsec:0604:02}

待撰写。

\subsection{群扩张与第二群上同调}
\label{b1:subsec:0604:03}

待撰写。


\begin{chaptersummary}
直积由进入各因子的同态刻画；交换群直和由从各因子出发的同态刻画。在一般群中，有限支集直积强制不同坐标交换，因而不具有群余积的泛性质。

\ProseParagraphBreak
半直积由两个群和一个自同构作用构成；指定扩张分裂，当且仅当其商映射具有同态截面。一般集合截面仍给出唯一坐标表达，但乘法中会出现因子集；反过来，满足相容恒等式的作用数据与因子集也能在 \(N\times H\) 上构造群扩张。更换截面只改变坐标数据，不改变原扩张；能否选择截面消去因子集，才决定扩张是否分裂。\(D_8\) 与 \(Q_8\) 的比较说明，即使核、商及其诱导作用相同，仍可能得到不同的扩张和不同的中间群。

\ProseParagraphBreak
扩张等价要求中间群的同构保持指定的核嵌入与商映射，比抽象群同构更严格。\chapref{b1:ch:07}将用正规列逐层分解群；本章说明了每一层的核与商怎样联系，以及为什么只知道各层因子通常不足以恢复整个群。
\end{chaptersummary}

\BookExercises

\begin{exercise}\label{b1:ex:06-1}
设 \(G\) 为群，\(\Delta=\bigl\{\,(g,g)\mid g\in G\,\bigr\}\leq G\times G\)。
\begin{enumerate}
\item 证明 \(\Delta\trianglelefteq G\times G\) 当且仅当 \(G\) 是交换群。
\item 设 \(G\neq\{1\}\)。证明 \(\Delta\) 的两个坐标投影均满射，但 \(\Delta\) 不是 \(G\times G\) 中两个坐标子群的直积。
\end{enumerate}
\end{exercise}
\begin{solution}
\textbf{（1）}
若 \(\Delta\) 正规，则对任意 \(a,g\in G\)，
\[
(a,1)(g,g)(a^{-1},1)=(aga^{-1},g)\in\Delta,
\]
故 \(aga^{-1}=g\)，所以 \(G\) 交换。反过来，若 \(G\) 交换，则 \(G\times G\) 交换，其中每个子群都正规，包括 \(\Delta\)。

\textbf{（2）}
当 \(G\neq\{1\}\) 时，\(\Delta\) 的两个坐标投影都是 \(G\)，但取 \(g\neq1\)，便有 \((g,1)\notin\Delta\)，所以 \(\Delta\neq G\times G\)。若 \(\Delta\) 是两个坐标子群的直积，这两个子群由投影都只能是 \(G\)，与上述严格不等矛盾。因此直积中的子群不必是两个因子中子群的坐标直积。
\end{solution}

\begin{exercise}\label{b1:ex:06-2}
设 \(m,n\) 为正整数。证明 \(C_m\times C_n\) 循环当且仅当 \(\gcd(m,n)=1\)，并求 \(C_4\times C_6\) 中元素阶的最大值。
\end{exercise}
\begin{solution}
若 \(a,b\) 分别生成 \(C_m,C_n\)，则 \((a,b)^d=1\) 当且仅当 \(m\mid d\) 且 \(n\mid d\)，故其阶为 \(\operatorname{lcm}(m,n)\)。任意元素的阶也整除这个最小公倍数。因此直积有阶为 \(mn\) 的生成元，当且仅当 \(\operatorname{lcm}(m,n)=mn\)，等价于 \(\gcd(m,n)=1\)。对 \(4,6\)，最大元素阶为 \(12\)，由两个生成元组成的元素达到，但直积阶为 \(24\)，所以不是循环群。
\end{solution}

\begin{exercise}\label{b1:ex:06-infinite-product-maps}
令 \(P=\prod_{r\geq1}\mathbb Z\)、\(S=\bigoplus_{r\geq1}\mathbb Z\)，并将 \(S\) 看成 \(P\) 的有限支集子群。记第 \(r\) 个坐标为 \(1\)、其他坐标为 \(0\) 的序列为 \(e_r\)。证明 \(S\) 由全部 \(e_r\) 生成，而 \(P\) 不由它们生成。再给出两个不同的同态 \(P\rightarrow P/S\)，使它们在每个坐标子群 \(\mathbb Z e_r\) 上的限制相同。
\end{exercise}
\begin{solution}
有限多个 \(e_r\) 的整数线性组合只有有限个非零坐标，故生成的子群包含于 \(S\)。反过来，若 \(a=(a_r)\in S\)，其支集 \(J\) 有限，就有 \(a=\sum_{r\in J}a_re_r\)。所以这些元素生成的子群恰为 \(S\)。常数序列 \(v=(1,1,\ldots)\) 不属于 \(S\)，因此它们不生成整个 \(P\)。

\(P\) 交换，故 \(S\trianglelefteq P\)，自然商映射 \(q:P\rightarrow P/S\) 与零同态 \(z:P\rightarrow P/S\) 都有定义。对每个 \(r\)，坐标子群 \(\mathbb Z e_r\) 包含于 \(S\)，所以两者在该子群上都为零。但 \(q(v)=v+S\neq S=z(v)\)，因而两个同态不同。从无限直积出发的同态，未必由它在各坐标子群上的限制决定。
\end{solution}

\begin{exercise}\label{b1:ex:06-3}
设 \(N\) 为加法交换群，\(H\) 为交换群，且 \(\alpha:H\rightarrow\operatorname{Aut}(N)\) 为群同态。令
\(N^H=\bigl\{\,n\in N\mid\alpha_h(n)=n\text{ 对所有 }h\in H\,\bigr\}\)。
证明 \(Z(N\rtimes_\alpha H)=N^H\times\ker\alpha\)，并计算非交换 \(21\) 阶群的中心。
\end{exercise}
\begin{solution}
元素 \((n,h)\) 与每个 \((m,1)\) 交换，当且仅当
\(n+\alpha_h(m)=m+n\) 对所有 \(m\) 成立，即 \(h\in\ker\alpha\)。它与每个 \((0,k)\) 交换，当且仅当 \(n=\alpha_k(n)\)，这里第二坐标交换用到了 \(H\) 交换。两种坐标子群生成整个半直积，所以这两个条件也充分。它们组成的群运算为直积运算，因为 \(\ker\alpha\) 作用平凡。对\cref{b1:thm:groups-order-pq}构造的非交换 \(21\) 阶群，\(N=\mathbb Z/7\mathbb Z\)、\(H=C_3\)，其生成元作用为乘以 \(2\)。共同固定元素必须满足 \(2n=n\)，所以只有 \(n=0\)；而三个作用分别为乘以 \(1,2,4\)，模 \(7\) 彼此不同，故 \(\ker\alpha=\{1\}\)。代入刚证明的中心公式，得到中心平凡。
\end{solution}

\begin{exercise}\label{b1:ex:06-4}
设整数 \(n\geq3\)。在 \(D_{2n}\) 中求旋转子群 \(\langle r\rangle\) 的全部补群。证明 \(n\) 为奇数时它们全共轭，\(n\) 为偶数时恰分成两个共轭类。
\end{exercise}
\begin{solution}
令 \(N=\langle r\rangle\)，\(\pi:D_{2n}\rightarrow D_{2n}/N\cong C_2\) 为商映射。若 \(H\) 是 \(N\) 的补群，则 \(NH=D_{2n}\) 说明 \(\pi|_H\) 满射，而 \(\ker(\pi|_H)=H\cap N=\{1\}\)。由\cref{b1:thm:first-isomorphism}，\(H\cong C_2\)，所以其非单位元素在旋转子群之外。\cref{b1:prop:dihedral-normal-form}说明所有这样的元素正是反射 \(r^is\)，且平方为单位。

反过来，每个 \(\langle r^is\rangle\) 与 \(N\) 交平凡，它在商群中的像含非单位元，故乘上 \(N\) 得到全群。因此全部补群为 \(\langle r^is\rangle\)，\(i\) 模 \(n\)，共有 \(n\) 个。

用 \(r^a\) 共轭使指数 \(i\) 变为 \(i+2a\)，用 \(s\) 共轭使指数变为 \(-i\)。若 \(n\) 奇，乘 \(2\) 在模 \(n\) 剩余类中可逆，任意两个指数都可由这种共轭互相变换。若 \(n\) 偶，这两种操作都保持指数奇偶；同奇偶指数又能通过加 \(2a\) 互相到达，因此恰有两个类。
\end{solution}

\begin{exercise}\label{b1:ex:06-5}
在实直线仿射群中，设 \(T\) 为平移子群。对 \(c\in\mathbb R\)，令 \(L_c\) 为固定点 \(c\) 的全部仿射变换。证明 \(L_c\) 是 \(T\) 的补群，并写出将 \(L_0\) 共轭到 \(L_c\) 的平移。
\end{exercise}
\begin{solution}
固定 \(c\) 的条件为 \(ac+b=c\)，故
\[
L_c=\Bigl\{\,\bigl((1-a)c,a\bigr)\mid a\in\mathbb R^\times\,\Bigr\}.
\]
它是点稳定子群，因而是子群；也可直接按半直积乘法验证。投影到 \(a\) 坐标在 \(L_c\) 上为同构，且 \(L_c\cap T=\{1\}\)。任意仿射元素有分解
\[
(b,a)=\bigl(b-(1-a)c,1\bigr)\bigl((1-a)c,a\bigr),
\]
所以 \(TL_c\) 为全群，\(L_c\) 是补群。令平移 \(\tau_c(x)=x+c\)，伸缩 \(f_a(x)=ax\)，则
\[
(\tau_c\circ f_a\circ\tau_c^{-1})(x)=a(x-c)+c.
\]
因此 \(\tau_c\) 把 \(L_0\) 共轭到 \(L_c\)。
\end{solution}

\begin{exercise}\label{b1:ex:06-complement-cocycle}
设 \(G=N\rtimes_\alpha H\)，并把 \(N\) 识别为 \(N\times\{1\}\)。对映射 \(b:H\rightarrow N\)，令
\[
L_b=\{\,\bigl(b(h),h\bigr)\mid h\in H\,\}.
\]
证明 \(L_b\) 是 \(N\) 在 \(G\) 中的补群，当且仅当
\[
b(1_H)=1_N,\qquad
b(hk)=b(h)\alpha_h\bigl(b(k)\bigr)
\quad(h,k\in H).
\]
并证明：用 \((n_0,1)\) 共轭 \(L_b\) 后所得补群对应的函数为
\[
b'(h)=n_0b(h)\alpha_h(n_0^{-1}).
\]
\end{exercise}
\begin{solution}
若题设两式成立，则
\[
\bigl(b(h),h\bigr)\bigl(b(k),k\bigr)
=\bigl(b(h)\alpha_h(b(k)),hk\bigr)
=\bigl(b(hk),hk\bigr).
\]
因此映射 \(h\mapsto\bigl(b(h),h\bigr)\) 是群同态，\(L_b\) 是其像，故为子群。第二坐标投影在 \(L_b\) 上与这个同态互逆，所以 \(L_b\cong H\) 且 \(L_b\cap N=\{1\}\)。任取 \((n,h)\in G\)，有
\[
(n,h)=\bigl(nb(h)^{-1},1\bigr)\bigl(b(h),h\bigr),
\]
故 \(NL_b=G\)，所以 \(L_b\) 是补群。

反过来，若 \(L_b\) 是补群，则第二坐标投影 \(L_b\rightarrow H\) 是同构：它由集合的定义满射，核为 \(L_b\cap N=\{1\}\)。其逆映射 \(h\mapsto\bigl(b(h),h\bigr)\) 必为同态；比较单位与乘积的第一坐标，就得到题设两式。

最后直接计算
\[
\begin{aligned}
(n_0,1)\bigl(b(h),h\bigr)(n_0^{-1},1)
&=\bigl(n_0b(h),h\bigr)(n_0^{-1},1)\\
&=\bigl(n_0b(h)\alpha_h(n_0^{-1}),h\bigr).
\end{aligned}
\]
因此共轭后的补群正是 \(L_{b'}\)，其中 \(b'\) 如题所述。这一区分了补群的存在、具体选择及不同补群是否共轭三个问题。
\end{solution}

\begin{exercise}\label{b1:ex:06-6}
设整数 \(m,n\geq2\)，采用 \(C_r=\mathbb Z/r\mathbb Z\) 的加法记号。考虑短正合列
\[
0\longrightarrow C_m\xrightarrow{\iota}C_{mn}\xrightarrow{\pi}C_n\longrightarrow0,
\]
其中 \(\iota\bigl([k]_m\bigr)=[nk]_{mn}\)，\(\pi\bigl([a]_{mn}\bigr)=[a]_n\)。证明这一扩张分裂，当且仅当 \(\gcd(m,n)=1\)。
\end{exercise}
\begin{solution}
先核对短正合性。改变 \(k\) 的代表会使 \(nk\) 改变 \(mn\) 的倍数，故 \(\iota\) 良定义；\(mn\mid nk\) 当且仅当 \(m\mid k\)，所以它单射。\(\pi\) 是满同态，其核恰由 \(n\) 的倍数的剩余类组成，也就是 \(\operatorname{im}\iota\)。

同态截面由生成元 \([1]_n\) 的像 \([a]_{mn}\) 确定，所需条件为
\[
a\equiv1\pmod n,\qquad na\equiv0\pmod{mn}.
\]
第二个条件等价于 \(m\mid a\)，写 \(a=mu\) 后，第一条件变为 \(mu\equiv1\pmod n\)。若有解，任何 \(m,n\) 的公因子都整除 \(mu-nv=1\)，所以二者互素；反过来，Bézout 等式在互素时给出这样的 \(u,v\)。取 \(a=mu\) 后，定义 \(s\Bigl([k]_n\Bigr)=[ka]_{mn}\)。若 \(k\) 改变 \(n\) 的倍数，\(ka\) 改变 \(na\) 的倍数，而 \(mn\mid na\)，所以该公式良定义；它保持加法，且 \(a\equiv1\pmod n\) 保证 \(\pi s\Bigl([k]_n\Bigr)=[k]_n\)。这就构造了同态截面。特别地，\(C_4\rightarrow C_2\) 不分裂，\(C_6\rightarrow C_3\) 分裂。
\end{solution}

\begin{exercise}\label{b1:ex:06-7}
设整数 \(n\geq2\)。对自然商映射 \(\mathbb Z\rightarrow\mathbb Z/n\mathbb Z\)，用 \(0,1,\ldots,n-1\) 作陪集代表，并通过 \(a\mapsto na\) 把核识别为 \(\mathbb Z\)。计算因子集，说明它怎样记录加法进位。
\end{exercise}
\begin{solution}
对整数 \(t\)，记 \(r_n(t)\in\{0,\ldots,n-1\}\) 为它除以 \(n\) 的余数；这是一个整数，与剩余类 \([t]_n\) 区分。对代表 \(0\leq i,j<n\)，有
\[
i+j=n\left\lfloor\frac{i+j}{n}\right\rfloor+r_n(i+j).
\]
核经 \(a\mapsto na\) 识别为 \(\mathbb Z\)，所以因子集为
\[
c(i,j)=\left\lfloor\frac{i+j}{n}\right\rfloor.
\]
其值为 \(0\) 或 \(1\)，记录一次加法是否产生进位。这里核和商群都采用加法记号，诱导作用平凡。\eqref{b1:eq:extension-factor-identity}具体成为
\[
\begin{aligned}
c(i,j)+c\Bigl(r_n(i+j),k\Bigr)
&=c(j,k)+c\Bigl(i,r_n(j+k)\Bigr)\\
&=\left\lfloor\frac{i+j+k}{n}\right\rfloor
\qquad(0\leq i,j,k<n).
\end{aligned}
\]
这表示先加哪两个数，都不改变最后的总进位。

在 \(\mathbb Z\times\{0,\ldots,n-1\}\) 上定义坐标运算
\[
(a,i)+(b,j)=\Bigl(a+b+c(i,j),r_n(i+j)\Bigr).
\]
整数除法说明 \((a,i)\mapsto na+i\) 是到 \(\mathbb Z\) 的双射；上面的余数等式又说明它把坐标运算变为整数加法。因此坐标运算确实给出一个群，并通过这个双射恢复原扩张。
\end{solution}

\begin{exercise}\label{b1:ex:06-8}
设 \(1\rightarrow A\rightarrow G\rightarrow H\rightarrow1\) 为中心扩张，\(H\) 为循环群。证明 \(G\) 交换，并说明这是否保证扩张分裂。另证明：任意满同态 \(\pi:G\rightarrow\mathbb Z\) 都有同态截面，这一结论不需要中心性假设。
\end{exercise}
\begin{solution}
取商群生成元的任一提升 \(x\in G\)。每个 \(g\in G\) 都可写为 \(ax^r\)，其中 \(a\in A\)、\(r\in\mathbb Z\)：商像写为生成元的幂后，将相应 \(x^r\) 除去，剩下的元素就在核中。因 \(A\leq Z(G)\)，对任意 \(a,b\in A\) 有
\[
(ax^r)(bx^s)=abx^{r+s}=bax^{s+r}=(bx^s)(ax^r).
\]
所以 \(G\) 交换。不过扩张不必分裂；\subsecref{b1:subsec:0603:02}已经验证模 \(2\) 的满射 \(C_4\rightarrow C_2\) 没有同态截面。这个例子表明，交换性只要求提升元素彼此交换，分裂还要求能选择满足有限循环商之阶数关系的提升。

对任意满同态 \(\pi:G\rightarrow\mathbb Z\)，取 \(x\in G\) 满足 \(\pi(x)=1\)，定义 \(s(k)=x^k\)。整数幂法则给出 \(s(k+\ell)=s(k)s(\ell)\)，而 \(\pi\bigl(s(k)\bigr)=k\)，所以 \(s\) 是同态截面。此处没有使用核的中心性。无限循环商没有有限阶关系需要提升元素满足；上面的不分裂例子来自有限循环商的这个额外要求。
\end{solution}

\begin{exercise}[综合题]\label{b1:ex:06-change-split-section}
设 \(C_4=\langle z\rangle\)、\(C_2=\{1,q\}\)，考虑第二坐标投影给出的扩张
\[
1\longrightarrow C_4\longrightarrow C_4\times C_2\longrightarrow C_2\longrightarrow1.
\]
\begin{enumerate}
\item 选择归一化集合截面 \(s(1)=(1,1)\)、\(s(q)=(z,q)\)。求其诱导作用和全部因子集值，并选取 \(b:C_2\rightarrow C_4\)，使新截面 \(s'(h)=(b(h),1)s(h)\) 的因子集恒为单位。解释这个例子为何不矛盾于分裂判别。
\item 令 \(N=S_3\)、\(H=C_2=\{1,q\}\)，考虑第二坐标投影 \(N\times H\rightarrow H\)。取 \(a=(1\,2\,3)\)，并选归一化集合截面
\[
s(1)=(1,1),\qquad s(q)=(a,q).
\]
计算相应的 \(\alpha,c\)，证明 \(\alpha:H\rightarrow\operatorname{Aut}(N)\) 不是群同态；再改变截面，使新的 \(\alpha'\) 平凡且 \(c'\equiv1\)。
\end{enumerate}
\end{exercise}
\begin{solution}
\textbf{（1）}
中间群 \(C_4\times C_2\) 交换，因此核上的诱导作用平凡。用第一坐标把核识别为 \(C_4\)，归一化给出
\[
c(1,1)=c(1,q)=c(q,1)=1,\qquad c(q,q)=z^2.
\]
最后一式来自 \(s(q)^2=(z^2,1)\)，故当前截面不是同态。

取 \(b(1)=1\)、\(b(q)=z^{-1}\)。将\eqref{b1:eq:factor-set-section-change}用于平凡作用，得到
\[
c'(h,k)=b(h)b(k)c(h,k)b(hk)^{-1}.
\]
有单位输入时，新因子集值仍为单位；唯一剩下的一对给出
\[
c'(q,q)=z^{-1}z^{-1}z^2=1.
\]
对应截面满足 \(s'(1)=(1,1)\)、\(s'(q)=(1,q)\)，正是同态截面。分裂要求存在这样的截面，并不要求所有集合截面都是同态；更换截面消去了这次坐标选择带来的非平凡因子集，扩张本身没有改变。

\textbf{（2）}
通过第一坐标把 \(N\) 识别为投影的核。截面对单位给出恒等自同构，而
\[
\alpha_q(n)=ana^{-1},
\qquad
c(1,1)=c(1,q)=c(q,1)=1,
\qquad
c(q,q)=a^2.
\]
最后一式来自 \(s(q)^2=(a^2,1)\)。这里
\[
\alpha_q^2=\operatorname{Inn}(a^2)\neq\operatorname{id}_N
=\alpha_{q^2};
\]
例如 \(\alpha_q^2\bigl((1\,2)\bigr)=(1\,3)\)。因此 \(\alpha\) 不是群同态。

取 \(b(1)=1\)、\(b(q)=a^{-1}\)，并令 \(s'(h)=(b(h),1)s(h)\)。于是
\[
s'(1)=(1,1),\qquad s'(q)=(1,q).
\]
新截面是第二坐标投影的标准同态截面；它的元素与 \(N\times\{1\}\) 逐元素交换，所以 \(\alpha'_q=\operatorname{id}_N\)，并且 \(s'(h)s'(k)=s'(hk)\) 给出 \(c'(h,k)=1\) 对所有 \(h,k\in H\) 成立。原扩张本来分裂，并不意味着任意集合截面产生的 \(\alpha\) 都是群作用。
\end{solution}
